\section{Discussion and limitations}
\label{sec:limits}
\paragraph{What the results say.}
On the two public frames, the investment needed to move a model to first place is governed jointly by the aggregation rule, the cap on per-task improvement and where the target stands. For a non-leading target the mean's cost is proportional to the gap ($L_1=n(\tau+\varepsilon-\bar x)$ when headroom suffices, Section~\ref{sec:setting}); rules that respond to the lowest or middle scores can be far cheaper in points, but a binding cap can make first place unreachable under them. The constructed examples show that the ranking of the mean, median and minimum rules can reverse, and the HELM census contradicts a blanket claim that the minimum must need the most tasks or points. Targeting matters less than an earlier version suggested (Section~\ref{sec:targeting}): a lowest-score-first greedy order attains the optimal task count at the ceiling for every score rule, so the optimiser adds value mainly in points and under a binding cap for the median rule.

\paragraph{Modelling assumptions.}
(i) The developer knows every score and all rivals are fixed; the minima hold only within this box model (fixed rivals, score-space cost),; public scores give no points-to-effort conversion. (ii) Tasks are controlled independently, although real training changes many tasks together. (iii) The caps are scenarios, not estimates of achievable gains. (iv) $L_1$ counts every point as equally costly. (v) ``First'' means strictly ahead by $0.001$ points; changing $\varepsilon$ to $0.1$ changed some task counts and point costs. (vi) The ordinal costs are exact only on the declared finite level set and upper-bound the real-valued infimum; infeasibility there is conditional on it; the Condorcet requirement is a sufficient-condition upper bound. (vii) The rival set is a convention (Section~\ref{sec:experiments}).

\paragraph{Data limits.}
Both frames are fixed censuses, not samples (HELM: 79 models, 57 correlated subjects; Open LLM Leaderboard: six tasks, merged and fine-tuned models): all ranges are sensitivity ranges, not confidence intervals.

\paragraph{Why the rules reverse between frames.}
The frames differ at once in task count, rivals, task correlation, score scale and population; matching, cap and unit change the picture (Section~\ref{sec:reversal}). We name no single cause and claim nothing for other leaderboards.

\paragraph{Reproduction.}
The original 16-file study and the review and final additions were each recomputed in isolated clean exports and agreed (maximum finite difference $5.8\times10^{-9}$ in the original, zero in the additions; Appendix~\ref{app:repro}). That pass used the bundled pinned processed matrices; re-downloading the raw sources was not repeated.
